\section{引言：课程概览}

本节课是 Stanford CS224N（Natural Language Processing with Deep Learning）的第三讲，由 Chris Manning 教授讲授。课程的核心目标是让学生理解神经网络训练的数学原理——如何\textbf{手动推导梯度}（矩阵微积分）以及如何\textbf{用算法自动计算梯度}（反向传播算法）。

\begin{figure}[H]
\centering
\includegraphics[width=0.95\textwidth]{slides-images/slide-000.jpg}
\caption{课程封面：CS224N Lecture 3——Neural net learning: Gradients by hand (matrix calculus) and algorithmically (the backpropagation algorithm)\protect\footnotemark}
\end{figure}
\footnotetext{Slides 第1页。}

课程分为三个部分：

\begin{figure}[H]
\centering
\includegraphics[width=0.95\textwidth]{slides-images/slide-009.jpg}
\caption{课程大纲：(1) 引言回顾（10分钟）；(2) 矩阵微积分（35分钟）；(3) 反向传播算法（35分钟）\protect\footnotemark}
\end{figure}
\footnotetext{Slides 第10页。}

\begin{importantbox}{本讲核心学习目标}
\begin{enumerate}
    \item 理解神经网络中梯度的数学推导——矩阵微积分（Matrix Calculus）
    \item 掌握反向传播算法（Backpropagation）——链式法则的高效实现
    \item 建立``前向传播计算函数值、反向传播计算梯度''的直觉
\end{enumerate}
\end{importantbox}

\subsection{本章小结}

本讲聚焦于神经网络训练的数学基础。Manning 教授强调，虽然现代深度学习框架（如 PyTorch）已经自动化了梯度计算，但理解底层原理对于调试、设计新模型和理解训练行为至关重要。

%% ========== Part 2: 神经网络基础回顾 ==========

